\chapter{About this blueprint}
\label{chap:intro}

This blueprint is a restructured, fine-grained version of the argument in
\texttt{Manuscript/eigenvalues.tex}. It serves two purposes:
\begin{itemize}
  \item a \emph{to-do list} for the Lean formalization, organized as a
        dependency graph;
  \item a place to \emph{digest} the manuscript: every node should be small
        enough that its statement is precise and its proof is a short,
        checkable argument.
\end{itemize}

\section*{How to read the nodes}

Every definition, lemma, proposition, theorem and corollary is a node of the
dependency graph. Remarks are not; we use them for commentary. Remarks titled
\emph{Digest note} record questions or problems found in the manuscript. They
are collected again in Chapter~\ref{chap:notes}.

Each node may carry the following annotations, which are invisible in the
pdf but drive the colours of the dependency graph in the web version.
\begin{itemize}
  \item \verb|\uses{...}| inside a statement: the definitions and results
        needed to \emph{state} it. Inside a proof: the results needed to
        \emph{prove} it.
  \item \verb|\lean{...}|: the name(s) of the corresponding Lean declaration(s).
  \item \verb|\leanok| inside a statement: the Lean statement exists (its
        proof may still be \texttt{sorry}). Inside a proof: the Lean proof is
        complete.
  \item \verb|\mathlibok|: the result is already in Mathlib.
  \item \verb|\notready|: the statement is not yet precise enough to be
        formalized; it needs to be digested first.
\end{itemize}
At the start nothing is formalized, so no node carries \verb|\lean| or
\verb|\leanok| yet.

\section*{Organization}

Chapters~\ref{chap:setup} and~\ref{chap:symmetry} set up the problem, state the
main theorem, and reduce it to four symmetry sectors.
Chapter~\ref{chap:geometry} collects the explicit geometry of the dipole flow
(identities for $\Psi$, the separatrix, the period function, orbit
comparisons). Chapter~\ref{chap:analysis} lists the analytic background that
the argument uses as a black box; this is the natural boundary between
``what we formalize'' and ``what we assume, at least initially''.
Chapters~\ref{chap:virial} and~\ref{chap:transport} isolate the two general
mechanisms used in every sector: virial (commutator) identities, and the
transport equation along orbits near the separatrix.
Chapters~\ref{chap:even-even}--\ref{chap:mixed} treat the sectors, with the
regularity theory needed for the odd sectors in Chapter~\ref{chap:regularity}.
Chapter~\ref{chap:certificates} collects the computer-assisted polynomial
inequalities.

\section*{Relation to the manuscript}

Manuscript labels are quoted in the form \texttt{(ee:trace-sos)} so they can
be found with a text search. The main structural changes compared to the
manuscript are:
\begin{itemize}
  \item Algebraic identities (commutators, sums of squares, coefficient
        computations) are separated from the analytic justification that
        they hold on the original $L^2$ domain. The algebraic nodes are the
        natural first formalization targets.
  \item The transport arguments that the manuscript reuses by referring back
        to the \emph{proofs} of earlier results (for instance
        ``apply the transport-only argument of
        Proposition~\texttt{prop:saddle-gains} to $(v,f)$'') are stated here as
        standalone lemmas about a general transport pair
        (Chapter~\ref{chap:transport}).
  \item Shared geometry of the flow is collected in one chapter instead of
        being derived inside sector-specific proofs.
\end{itemize}
